\documentclass{article}
\usepackage[utf8]{inputenc}
\usepackage{listings}
\usepackage{xcolor}
\usepackage{geometry}
\usepackage{graphicx}
\usepackage{fancyhdr}
\usepackage{sectsty}
\usepackage{enumitem}
\usepackage{hyperref}

\geometry{margin=1in}

% Colors for code
\definecolor{codebg}{rgb}{0.95,0.95,0.95}
\definecolor{keyword}{rgb}{0.36,0.09,0.54}
\definecolor{comment}{rgb}{0.0,0.5,0.0}
\definecolor{string}{rgb}{0.55,0.0,0.0}

% Listings style
\lstset{
    backgroundcolor=\color{codebg},
    basicstyle=\ttfamily\footnotesize,
    keywordstyle=\color{keyword}\bfseries,
    commentstyle=\color{comment}\itshape,
    stringstyle=\color{string},
    numbers=left,
    numberstyle=\tiny,
    stepnumber=1,
    numbersep=5pt,
    tabsize=4,
    showstringspaces=false,
    breaklines=true
}

% Header and footer
\pagestyle{fancy}
\fancyhf{}
\lhead{Department of Computer Science and Engineering}
\rhead{University of Dhaka}
\cfoot{\thepage}

\sectionfont{\centering \Large \bfseries}
\subsectionfont{\centering \normalsize \bfseries}

% Title page content
\begin{document}

\begin{titlepage}
    \centering
    \vspace*{1cm}
    \includegraphics[width=0.3\textwidth]{csedu.png} \\[1cm]
    
    {\LARGE Department of \\ Computer Science and Engineering \\[0.5cm]}
    {\Large University of Dhaka \\[2cm]}
    
    \hrule
    \vspace{1cm}
    {\LARGE Title: Implementation of Client-Server \\ Communication using Socket Programming \\[1cm]}
    \hrule
    \vspace{1cm}
    
    {\large CSE 3111: Computer Networking Lab \\[0.5cm]}
    {\normalsize Batch: 29/3rd Year 1st Semester 2024\\[0.5cm]}
    
    {\large Prepared by: Mehedi Hasan (22) and Biplob Pal (50)}
    
    \vfill
\end{titlepage}

\section{Introduction}
Socket programming allows communication between two programs over a network. Each program communicates through endpoints called sockets. Basic socket programming allows only a single client to connect at a time, which can create delays and reduce efficiency for multiple users.

Multi-threaded socket programming enables a server to handle multiple clients simultaneously by assigning each client connection to a separate thread. This approach is essential for real-time applications, such as chat or file-sharing systems, where multiple users must interact concurrently without blocking each other.

\section{Objectives}
The main objectives of this lab are:
\begin{enumerate}
    \item To implement a multi-threaded HTTP file server and client using Java.
    \item To learn file transfer over a network using GET and POST methods.
    \item To understand the advantages of multi-threaded socket programming over single-threaded servers.
\end{enumerate}

\section{Design Details}
The implementation uses a client-server architecture:

\subsection{Server Design Steps}
\begin{enumerate}
    \item Initialize a server socket on port 5000.
    \item Create a folder to store files (``server\_files'').
    \item Accept incoming client connections.
    \item For each client, start a new thread to handle requests.
    \item Supported client commands:
    \begin{itemize}
        \item \texttt{ls} -- list all files in the server folder.
        \item \texttt{filename.ext} -- download a specific file.
        \item \texttt{exit} -- close connection.
    \end{itemize}
\end{enumerate}

\subsection{Client Design Steps}
\begin{enumerate}
    \item Connect to the server using IP and port.
    \item Provide a command-line interface to:
    \begin{itemize}
        \item List available files on the server.
        \item Download files using GET requests.
        \item Upload files using POST requests.
        \item Exit the session.
    \end{itemize}
    \item Receive file data in chunks and save locally.
\end{enumerate}

\subsection{Operation Flow (Textual Representation)}
\begin{enumerate}
    \item Server starts and waits for client connections.
    \item Client connects and receives a welcome message.
    \item Client sends a command:
    \begin{itemize}
        \item If \texttt{ls}, server sends the list of files.
        \item If \texttt{filename.ext}, server checks existence and sends file.
        \item If \texttt{exit}, server closes the connection.
    \end{itemize}
    \item Client receives response and saves files if requested.
\end{enumerate}

\section{Implementation}
\subsection{Server Code}
\begin{lstlisting}[language=Java]
import java.io.*;
import java.net.*;

public class Roll50_Server {
    private static final int PORT = 5000;
    private static final String Folder = "server_files";

    public static void main(String[] args) {
        File folder = new File(Folder);
        if (!folder.exists()) folder.mkdir();

        try (ServerSocket serverSocket = new ServerSocket(PORT)) {
            System.out.println("Server started. Listening on port " + PORT);
            while (true) {
                Socket socket = serverSocket.accept();
                System.out.println("Client connected: " + socket.getInetAddress());
                new ClientHandler(socket, folder).start();
            }
        } catch (IOException e) { e.printStackTrace(); }
    }
}

class ClientHandler extends Thread {
    private Socket socket;
    private File folder;

    public ClientHandler(Socket socket, File folder) {
        this.socket = socket;
        this.folder = folder;
    }

    @Override
    public void run() {
        try (BufferedReader reader = new BufferedReader(new InputStreamReader(socket.getInputStream()));
             PrintWriter writer = new PrintWriter(socket.getOutputStream(), true);
             DataOutputStream dataOut = new DataOutputStream(socket.getOutputStream())) {

            writer.println("Connected to File Server. Type 'ls' to see files or enter a file name.");

            String command;
            while ((command = reader.readLine()) != null) {
                if (command.equalsIgnoreCase("ls")) {
                    File[] files = folder.listFiles();
                    if (files != null) for (File f : files) if (f.isFile()) writer.println(f.getName());
                    writer.println("END");
                } else {
                    File file = new File(folder, command);
                    if (file.exists() && file.isFile()) {
                        writer.println("FOUND");
                        long fileSize = file.length();
                        dataOut.writeLong(fileSize);
                        try (FileInputStream fis = new FileInputStream(file)) {
                            byte[] buffer = new byte[4096];
                            int bytesRead;
                            while ((bytesRead = fis.read(buffer)) != -1) {
                                dataOut.write(buffer, 0, bytesRead);
                            }
                        }
                        System.out.println("File " + command + " sent successfully.");
                    } else {
                        writer.println("NOT_FOUND");
                        System.out.println("File " + command + " not found.");
                    }
                }
            }
        } catch (IOException e) { e.printStackTrace(); } 
        finally { try { socket.close(); } catch (IOException e) { e.printStackTrace(); } }
    }
}
\end{lstlisting}

\subsection{Client Code}
\begin{lstlisting}[language=Java]
import java.io.*;
import java.net.*;

public class Roll22_Client {
    static final String SERVER_IP = "10.101.198.21";
    static final int SERVER_PORT = 5000;

    public static void main(String[] args) {
        try (Socket socket = new Socket(SERVER_IP, SERVER_PORT);
             BufferedReader reader = new BufferedReader(new InputStreamReader(socket.getInputStream()));
             PrintWriter writer = new PrintWriter(socket.getOutputStream(), true);
             DataInputStream dtinp = new DataInputStream(socket.getInputStream());
             BufferedReader bfr = new BufferedReader(new InputStreamReader(System.in))) {

            System.out.println("Connected to server!");
            System.out.println(reader.readLine());

            String command;
            while (true) {
                System.out.print("Enter command (ls / filename.ext / exit): ");
                command = bfr.readLine();
                if (command.equalsIgnoreCase("exit")) break;

                writer.println(command);

                if (command.equalsIgnoreCase("ls")) {
                    String fileName;
                    while (!(fileName = reader.readLine()).equals("END")) {
                        System.out.println(fileName);
                    }
                } else {
                    String serverResponse = reader.readLine();
                    if (serverResponse.equals("FOUND")) {
                        long fileSize = dtinp.readLong();
                        FileOutputStream fos = new FileOutputStream("downloaded_" + command);
                        byte[] buffer = new byte[4096];
                        int bytesRead;
                        long totalRead = 0;
                        while (totalRead < fileSize && (bytesRead = dtinp.read(buffer)) != -1) {
                            fos.write(buffer, 0, bytesRead);
                            totalRead += bytesRead;
                        }
                        fos.close();
                        System.out.println("Downloaded done: downloaded_" + command);
                    } else {
                        System.out.println("File not found on server!");
                    }
                }
            }
        } catch (Exception e) { e.printStackTrace(); }
    }
}
\end{lstlisting}

\section{Result Analysis}
\subsection{Server Output (Textual)}
\begin{verbatim}
Server started. Listening on port 5000
Client connected: /192.168.0.101
File example.txt sent successfully.
File not found.
\end{verbatim}

\subsection{Client Output (Textual)}
\begin{verbatim}
Connected to server!
Connected to File Server. Type 'ls' to see files or enter a file name.
Enter command: ls
example.txt
notes.pdf
END
Enter command: example.txt
Downloaded done: downloaded_example.txt
Enter command: unknown.txt
File not found on server!
Enter command: exit
\end{verbatim}

\subsection{Description}
The server handles multiple clients and responds correctly to commands. The client can list files, download files, and is notified if a file does not exist.

\section{Discussion}
Basic socket programming only supports one client at a time. This causes:
\begin{itemize}
    \item Delays for other clients.
    \item Poor performance for multiple users.
    \item Unresponsive server in real-time applications.
\end{itemize}

Multi-threaded socket programming solves these issues by creating a separate thread for each client:
\begin{itemize}
    \item Multiple clients can connect simultaneously.
    \item Server is responsive and scalable.
    \item Supports concurrent file transfers.
\end{itemize}

During this lab, I learned how to implement a multi-threaded server, manage concurrent connections, and handle file transfer efficiently. Challenges included handling client disconnections and ensuring proper file stream management.

\section{Conclusion}
The lab demonstrates a multi-threaded HTTP-like file server and client in Java. It shows the advantages over basic socket programming and provides hands-on experience with network programming and concurrent file transfer.

\end{document}
